\documentclass[UTF8,a4paper,11pt]{ctexart}

\usepackage[margin=2.4cm]{geometry}
\usepackage{amsmath,amssymb}
\usepackage{booktabs}
\usepackage{array}
\usepackage{enumitem}
\usepackage{xcolor}
\usepackage{titlesec}
\usepackage{fancyhdr}
\usepackage{caption}
\usepackage[colorlinks=true,linkcolor=blue!60!black,citecolor=blue!60!black,urlcolor=blue!60!black]{hyperref}

\definecolor{mainblue}{RGB}{20,60,120}
\definecolor{accent}{RGB}{180,70,30}

\titleformat{\section}{\Large\bfseries\color{mainblue}}{\thesection}{0.6em}{}
\titleformat{\subsection}{\large\bfseries\color{accent}}{\thesubsection}{0.6em}{}
\captionsetup{font=small,labelfont=bf}
\setlist[itemize]{leftmargin=1.6em,itemsep=2pt,topsep=3pt}
\setlist[enumerate]{leftmargin=1.8em,itemsep=2pt,topsep=3pt}

\pagestyle{fancy}
\fancyhf{}
\fancyhead[L]{\small SFSE 阶段五技术说明}
\fancyhead[R]{\small\thepage}
\renewcommand{\headrulewidth}{0.4pt}

\title{\vspace{-1.4cm}\textbf{\color{mainblue}安全可行主从博弈均衡\\阶段五技术说明与项目总结}\\[6pt]
\large 全尺寸安全 RL、不确定性与极端天气鲁棒性}
\author{充电站微网数字孪生项目组}
\date{\today}

\begin{document}
\sloppy
\maketitle
\thispagestyle{fancy}

\begin{center}
\begin{minipage}{0.93\textwidth}
\small
\textbf{摘要：}本说明报告 SFSE 研究第五阶段成果并总结全项目。阶段五在 24 时段全尺寸算例上，用多起点 BRI 均衡解作为示范标签，训练全尺寸安全 RL 策略；并评估不确定性与极端天气下的鲁棒性。结果：全尺寸安全 RL 相对 BRI 的净利润间隙为 $-13.05\%$，但\textbf{定价决策仅需 0.69 ms}（相对 BRI 约 $4\times10^4$ 倍加速）；在 20 个不确定性场景下，SFSE-BRI 的平均净利润 2043.7 元、最差 $-161.7$ 元，均优于固定电价的 1112.0 元与 $-1800.5$ 元；极端天气下 SFSE 的失负荷（17068/17411 kWh）低于固定电价（19168/19611 kWh）。
\end{minipage}
\end{center}

\vspace{0.4cm}

\section{阶段五计划}
\begin{enumerate}
  \item 全尺寸安全 RL：以 BRI 均衡解为示范，行为克隆预训练 + A2C 微调；
  \item 不确定性鲁棒性：风光/负荷/用户扰动下的收益分布与可靠性；
  \item 极端天气鲁棒性：新能源骤降时的失负荷对比；
  \item 项目总结与展望。
\end{enumerate}

\section{全尺寸安全强化学习}

\subsection{方法}
在全尺寸算例（$T=24$、8 用户群）上采样 20 个随机场景，用多起点 BRI 求示范均衡（共耗时 527 s）；以场景特征为输入、BRI 电价为标签，行为克隆预训练策略均值，再用 A2C 微调 800 episodes。

\subsection{结果}
\begin{table}[htbp]
\centering
\small
\caption{全尺寸安全 RL vs BRI（4 个测试场景）}
\begin{tabular}{p{4.4cm} p{4.4cm} p{4.4cm}}
\toprule
\textbf{指标} & \textbf{安全 RL} & \textbf{BRI 均衡}\\
\midrule
平均净利润 & 1735.2 元 & 1995.5 元\\
最优性间隙 & \multicolumn{2}{c}{$-13.05\%$}\\
定价决策时间 & \textbf{0.69 ms} & 27291 ms\\
决策加速比 & \multicolumn{2}{c}{$\mathbf{3.97\times10^4}$ 倍}\\
可靠性（安全层） & 100\% & 100\%\\
\bottomrule
\end{tabular}
\end{table}

\section{不确定性与极端天气鲁棒性}

\subsection{不确定性场景（20 个）}
\begin{table}[htbp]
\centering
\small
\caption{不确定性场景下的鲁棒性对比}
\begin{tabular}{p{4.4cm} p{4.6cm} p{4.6cm}}
\toprule
\textbf{指标} & \textbf{SFSE-BRI} & \textbf{固定电价 1.2}\\
\midrule
平均净利润 & \textbf{2043.7 元} & 1112.0 元\\
最差净利润 & \textbf{$-161.7$ 元} & $-1800.5$ 元\\
可靠率（无失负荷） & 65\% & 70\%\\
\bottomrule
\end{tabular}
\end{table}

SFSE 在平均收益（$+84\%$）与最差情形上均显著优于固定电价；固定电价可靠率略高，是因为其低电价诱导更多负荷、但收益低，属"低价换可靠"的权衡。

\subsection{极端天气（新能源骤降 50\%--75\%）}
\begin{table}[htbp]
\centering
\small
\caption{极端天气失负荷对比（kWh）}
\begin{tabular}{p{4.4cm} p{4.4cm} p{4.4cm}}
\toprule
\textbf{场景} & \textbf{SFSE-BRI} & \textbf{固定电价 1.2}\\
\midrule
极端场景 1 & \textbf{17068} & 19168\\
极端场景 2 & \textbf{17411} & 19611\\
\bottomrule
\end{tabular}
\end{table}

极端天气下两者都会失负荷（离网系统物理受限），但 SFSE 的失负荷更低，说明博弈定价能通过引导负荷错峰、配合储能减缓缺额。

\section{五阶段项目总结}

\begin{table}[htbp]
\centering
\small
\caption{五阶段成果总览}
\begin{tabular}{p{1.2cm} p{4.4cm} p{5.0cm} p{3.0cm}}
\toprule
\textbf{阶段} & \textbf{主题} & \textbf{关键结果} & \textbf{产出}\\
\midrule
一 & 小算例 KKT-MILP & 均衡真解 415.66 元 & run\_small\_case\\
二 & BRI + 安全层 + 命题 & 可行率/存在性 100\%，间隙 0 & Stage2 说明\\
三 & 代理 + 安全 RL & 代理 425$\times$，RL 间隙 $-1.24\%$ & Stage3 说明\\
四 & 全尺寸基准 + SOH & 收益 $+37.8\%$，SOH 更低 & Stage4 说明\\
五 & 全尺寸 RL + 鲁棒性 & 决策 $4\times10^4\times$，极端更优 & 本说明\\
\bottomrule
\end{tabular}
\end{table}

\subsection{核心贡献}
\begin{enumerate}
  \item 提出\textbf{安全可行主从博弈均衡（SFSE）}新问题，将离网物理约束内嵌于均衡可行域；
  \item 建立\textbf{定价—调度—寿命}三角耦合模型，量化收益与 SOH 的协同；
  \item 设计\textbf{结构化安全投影层}与\textbf{可微最佳反应代理}，保证可行并实现数量级加速；
  \item 给出均衡\textbf{存在性、安全性、收敛性}命题与数值验证；
  \item 揭示精确 MILP 的\textbf{全尺寸规模爆炸}，并给出 BRI/RL 可扩展求解路线。
\end{enumerate}

\subsection{局限与展望}
\begin{itemize}
  \item 全尺寸安全 RL 的间隙（$-13\%$）仍偏大，需更多示范数据或更强的序列建模；
  \item 随机场景的可靠率偏低，源于离网系统在低新能源时物理受限，可引入需求响应与备用容量；
  \item 电池退化采用线性化吞吐 + 日历近似，未来可嵌入完整非线性 SOH 模型；
  \item 下一步：接入真实交通流与数字孪生在线数据，形成闭环在线学习与工程化部署。
\end{itemize}

\end{document}
